%% kapa: exam=1 round=26 year=2015 period=1 form=A subject=civil_law subject_ko=민법 seq=25 number=25
%% source: Q-Net 감정평가사 기출문제 (cst00309) artlSeq=5138001 file=extracted/1008/5138001/제26회_감정평가사 1차_기출/제26회_감정평가사 1차_기출/1교시_감평_A형.pdf
%% conversion: kordoc 4.24.51 -> tools/gen1.py (auto); review=proofread-sonnet; flags=none
\begin{question}{25}
중간생략등기에 관한 설명으로 옳지 않은 것은? (다툼이 있으면 판례에 따름)
\begin{choices}
  \item 중간생략등기의 합의는 순차적 또는 묵시적으로 할 수 있다.
  \item 중간생략등기의 합의가 있더라도 최초매도인과 최종매수인 사이에 매매계약이 체결되었다고 볼 수는 없다.
  \item 중간생략등기의 합의가 있다고 하여 최초매도인이 매매계약상 상대방에 대하여 가지는 대금청구권의 행사가 제한되는 것은 아니다.
  \item 관계당사자 전원의 의사합치가 없어도 중간자의 동의가 있다면 최종매수인은 최초매도인을 상대로 직접 중간생략등기를 청구할 수 있다.
  \item 중간생략등기가 당사자 사이에 적법한 등기원인에 기하여 이미 경료되었다면, 중간생략등기의 합의가 없었음을 들어 그 등기의 말소를 구할 수는 없다.
\end{choices}
\end{question}
\answer{4}
